\subsubsection{Goldstoneボソンの作用}

ユニタリーゲージのラグランジアンは非常に一般的であるが、明らかにあまり直観的ではない。例えば、ある特定の極限では標準的なスローロール・インフレーションを含む。しかし、スカラーの自由度はどこにあるのだろうか。これほど複雑なのは、これが自発的に破れたゲージ対称性のユニタリーゲージでのラグランジアンだからである。

\noindent\textbf{Goldstoneボソンの等価定理：}
ユニタリーゲージのラグランジアンは、重力子の二つのヘリシティーと一つのスカラーモードという三つの自由度を記述する。このモードは、破れた時間の微分同相変換を行うことにより（Stückelbergの手法）、この対称性を非線形に実現するGoldstoneボソンとして明示的に現れる。質量をもつゲージボソンの縦成分に関する等価定理\cite{Cornwall:1974km}との類推から、混合を無視できる短距離では、Goldstoneボソンの物理は重力子の二つのヘリシティーから分離すると期待される。同じ方法をわれわれの場合に適用する前に、非Abelianゲージ理論で何が起こるかを簡単に復習しよう。

非Abelianゲージ群の$A_\mu^a$に対するユニタリーゲージの作用は
\OriginalEquation{166}
である。ここで$A_\mu=A_\mu^aT^a$である。ゲージ変換の下で
\OriginalEquation{167}
となる。したがって作用は
\OriginalEquation{168}
となる。質量項はゲージ不変ではなかったため、この項には$U$の因子が現れる。$U=\exp[iT^a\pi^a(t,\vec{x})]$と書くことにより、ゲージ不変性を「回復」できる。ここで$\pi^a$はスカラー（Goldstoneボソン）であり、ゲージ変換$\Lambda$の下で
\OriginalEquation{169}
のように非線形に変換する。

ここで一時的に、ゲージ理論が$U(1)$理論である場合を考えると、
\OriginalEquation{170}
となることに注意しよう。$\pi$はゲージ変換の下でシフトする。$0$が$0$へ写されないので、これは非線形変換である。非線形に変換する力学的な場を再導入することによって、ゲージ不変性が回復された。ゲージ不変性は非線形に実現されている。

正準規格化に移ると、
\OriginalEquation{171}
Goldstoneボソンの自己相互作用は、ゲージボソンの質量よりもパラメトリックに高いスケール$4\pi m/g$で強結合になることが分かる。Goldstoneボソンを再導入する利点は、エネルギー$E\gg m$では、Goldstoneボソンとゲージ場の横成分との混合が重要でなくなり、両者のセクターが分離することである。実際、式\eqref{eq:167}の混合項は
\OriginalEquation{172}
という形であり、$E\gg m$では正準運動項$(\partial\pi_c)^2$に比べて重要でない。

$D_\mu U D^\mu U$という項を展開すると、
\OriginalEquation{173}
という形の非関連な（すなわち繰り込み不可能な）項が得られることに注意しよう。これは、エネルギー$E\sim4\pi m/g$で強結合になり、ユニタリー性の破れを引き起こす演算子である。

$m\ll E\ll4\pi m/g$の範囲では、Goldstoneボソン$\pi$の物理は弱結合であり、横成分との混合を無視して調べることができる。

われわれが扱う、時間の微分同相変換が破れた場合について、同じ手順をたどろう。例えば、次の二つの演算子に注目しよう。
\OriginalEquation{174}
破れた時間の微分同相変換$t\to\tilde{t}=t+\xi^0(x)$、$\vec{x}\to\tilde{\vec{x}}=\vec{x}$の下で、$g^{00}$は次のように変換する。
\OriginalEquation{175}
変換後の場を用いて書いた作用は、次のようになる。
\OriginalEquation{176}
積分変数を$\tilde{x}$に変えると、次を得る。
\OriginalEquation{177}
Goldstoneボソンを再導入する手順は、ここでもゲージ理論の場合と同様である。上の作用に$\xi^0$が現れる箇所ではすべて、次の置き換えを行う。
\OriginalEquation{178}
